\documentclass[11pt,a4paper]{article}
\usepackage[T1]{fontenc}
\usepackage[utf8]{inputenc}
\usepackage{amsmath,amssymb,xcolor}
\usepackage[margin=25mm]{geometry}
\usepackage[colorlinks=true,urlcolor=blue]{hyperref}
\definecolor{nexorablue}{HTML}{2459B5}
\setlength{\parindent}{0pt}
\setlength{\parskip}{7pt}
\setlength{\emergencystretch}{2em}
\renewcommand{\familydefault}{\sfdefault}
\begin{document}

{\small\bfseries\color{nexorablue} NEXORA CPGE}\par
{\small MP / NEX-IND-05 / Corrigé}\par
{\LARGE\bfseries Courant transitoire dans une bobine secondaire}\par
\vspace{0.25em}\hrule\vspace{0.6em}

\section*{Équation de la boucle}
L’équation s’écrit $L\dot i_2+Ri_2=-M\dot i_1=-0{,}20\,\mathrm V$. La constante de temps est $0{,}040\,\mathrm s$.

\section*{Solution}
La limite pour une rampe prolongée est $-0{,}20/3\,\mathrm A$. Ainsi $i_2(t)=-(0{,}20/3)(1-e^{-t/0{,}040})\,\mathrm A$. À $80\,\mathrm{ms}$, le résultat vaut environ $-0{,}0576\,\mathrm A$.


\end{document}
